\chapter{Apresentação}

\section*{\textcolor{azulPrincipal}{\faRocket} Bem-vindo(a) ao mundo da programação!}

\begin{destacado}
\textit{Se você está lendo esta apostila, é porque está dando o primeiro passo em uma jornada muito legal: aprender a programar usando a \textbf{linguagem C}. Não se preocupe se você nunca escreveu uma linha de código antes --- todo programador, por mais experiente que seja, começou exatamente como você está começando agora.}
\end{destacado}

\vspace{0.5cm}

O objetivo deste capítulo é simples: explicar, de forma bem tranquila, \textbf{o que é programar}, um pouco da \textbf{história da computação e da linguagem C}, \textbf{para que ela é usada até hoje}, \textbf{qual é o seu legado}, \textbf{como um programa em C é organizado} e \textbf{onde vamos escrever e executar nossos códigos}. Vamos usar comparações do dia a dia para facilitar o entendimento.

\section{\textcolor{azulPrincipal}{O que é programar?}}

Programar é \textbf{\textcolor{laranja}{dar instruções para o computador executar}}.

\begin{importante}[Pense nisto]
Pense no computador como uma pessoa extremamente obediente, mas que só entende ordens muito específicas, na ordem certa, sem pular nenhum passo. Se você pedir para essa pessoa ``fazer um suco'', ela não vai entender. Mas se você disser:
\end{importante}

\begin{enumerate}
    \item Pegue duas laranjas.
    \item Corte as laranjas ao meio.
    \item Esprema as laranjas no espremedor.
    \item Coloque o suco em um copo.
\end{enumerate}

Aí sim! Cada passo é uma \textbf{instrução}, e o conjunto de instruções, na ordem certa, é o que chamamos de \textbf{programa} (ou \textbf{algoritmo}, quando ainda não está em uma linguagem específica).

A \textbf{linguagem de programação} é justamente o ``idioma'' que usamos para escrever essas instruções de um jeito que o computador consiga entender e executar.

\section{\textcolor{azulPrincipal}{Uma pequena viagem pela história da computação}}

Para entender por que aprendemos C hoje em dia, é legal entender de onde tudo veio.

\begin{itemize}
    \item \textcolor{laranja}{\textbf{Décadas de 1940 e 1950}}: surgem os primeiros computadores eletrônicos, enormes (do tamanho de uma sala!) e que precisavam ser programados quase que ``ligando fios'' e usando código de máquina (sequências de 0s e 1s). Era um trabalho extremamente difícil e demorado.
    
    \item \textcolor{laranja}{\textbf{Final dos anos 1950 e década de 1960}}: surgem as primeiras linguagens de ``alto nível'', como o Fortran e o COBOL, que tornaram a programação um pouco mais parecida com a linguagem humana, e menos parecida com ``fios e números binários''.
    
    \item \textcolor{laranja}{\textbf{Final dos anos 1960}}: pesquisadores perceberam que era preciso de uma linguagem que fosse, ao mesmo tempo, \textbf{fácil de escrever} para humanos e \textbf{eficiente} o suficiente para controlar diretamente o hardware do computador. Foi nesse cenário que a linguagem C começou a ser criada.
\end{itemize}

Ou seja: o C nasceu da necessidade de ter uma linguagem que unisse \textbf{\textcolor{verdeC}{praticidade para o programador}} com \textbf{\textcolor{verdeC}{controle e desempenho próximos ao do hardware}}.

\section{\textcolor{azulPrincipal}{O surgimento da linguagem C}}

A linguagem C foi criada por \textbf{\textcolor{laranja}{Dennis Ritchie}}, nos \textbf{Laboratórios Bell} (Bell Labs), nos Estados Unidos, durante o início da década de \textbf{1970} (entre 1969 e 1973, aproximadamente).

O motivo principal foi bem prático: a equipe estava desenvolvendo um sistema operacional chamado \textbf{Unix}, e precisava de uma linguagem que permitisse:

\begin{destacado}
\begin{itemize}
    \item Escrever um sistema operacional inteiro de forma mais organizada (até então, sistemas operacionais eram escritos quase totalmente em \textbf{Assembly}, uma linguagem muito próxima do código de máquina e muito difícil de manter).
    \item Ter \textbf{bom desempenho}, já que estava sendo usada para controlar diretamente o computador.
    \item Ser \textbf{portátil}, ou seja, possível de usar em diferentes tipos de computador, com poucas alterações.
\end{itemize}
\end{destacado}

O resultado foi tão bom que, pouco tempo depois, \textbf{grande parte do próprio sistema Unix foi reescrita em C}. Esse foi um marco enorme: mostrou que era possível construir até mesmo um sistema operacional com uma linguagem relativamente simples e organizada, e não apenas em Assembly.

\section{\textcolor{azulPrincipal}{Para que a linguagem C é usada hoje em dia?}}

Mesmo tendo sido criada há mais de 50 anos, o C continua extremamente utilizado. Alguns exemplos de onde ele aparece:

\begin{dica}[O C está em toda parte]
\begin{itemize}
    \item \textbf{Sistemas operacionais}: partes do núcleo (kernel) do Linux, do Windows e de outros sistemas são escritas em C.
    \item \textbf{Sistemas embarcados}: programas que rodam dentro de eletrodomésticos, sensores, equipamentos industriais, dispositivos médicos, controladores de carros, entre outros.
    \item \textbf{Drivers de dispositivos}: os programas que fazem o sistema operacional ``conversar'' com impressoras, placas de vídeo, placas de rede, etc.
    \item \textbf{Linguagens e ferramentas}: muitos compiladores, interpretadores de outras linguagens (como o próprio Python) e bancos de dados (como o MySQL e o PostgreSQL) têm partes importantes escritas em C.
    \item \textbf{Jogos e engines gráficas}: por seu alto desempenho, o C (e o C++, que vem dele) ainda é usado em partes críticas de jogos e motores gráficos.
\end{itemize}
\end{dica}

Ou seja: sempre que você usa um computador, celular, smart TV, carro moderno ou eletrodoméstico ``inteligente'', existe uma boa chance de que, em algum lugar ali dentro, exista código em C rodando.

\section{\textcolor{azulPrincipal}{O legado da linguagem C}}

O \textbf{legado} do C vai muito além do código que ainda roda por aí. Ele influenciou diretamente a forma como \textbf{diversas outras linguagens foram criadas}. Algumas delas:

\begin{itemize}
    \item \textbf{C++}: criado como uma extensão do C, adicionando recursos de orientação a objetos.
    \item \textbf{Java}: pegou muito da sintaxe do C (uso de chaves \texttt{\{\}}, ponto e vírgula \texttt{;}, estrutura de funções, etc.).
    \item \textbf{C\#}: criado pela Microsoft, também segue uma sintaxe muito parecida com a do C.
    \item \textbf{JavaScript}, \textbf{PHP}, \textbf{Python}, entre outras: também foram influenciadas, em maior ou menor grau, pela forma como o C organiza seu código.
\end{itemize}

Por isso, quando dizemos que o C é a ``língua-mãe'' de várias linguagens modernas, não é exagero. Aprender C é, de certa forma, aprender a \textbf{base de um ``idioma de programação''} que se espalhou pelo mundo todo, ganhando ``dialetos'' (outras linguagens) ao longo dos anos.

Além disso, muitos conceitos fundamentais de programação que usamos até hoje --- como variáveis, estruturas de controle (\texttt{if}, \texttt{for}, \texttt{while}), funções e ponteiros --- foram popularizados justamente através do C, e aparecem (com pequenas variações) em praticamente todas as linguagens modernas.

\section{\textcolor{azulPrincipal}{Curiosidade: o ``Olá, Mundo!''}}

Você provavelmente vai escrever, em algum momento, um programa que apenas exibe a frase \textbf{``Hello, World!''} (ou ``Olá, Mundo!'') na tela. Isso não é coincidência!

\begin{atencao}[Curiosidade]
Esse exemplo ficou famoso justamente por causa de um livro sobre a linguagem C, escrito pelo próprio criador da linguagem, \textbf{Dennis Ritchie}, em conjunto com \textbf{Brian Kernighan}, chamado \textit{``The C Programming Language''} (lançado em 1978). Esse livro usava o programa ``Hello, World!'' como o primeiro exemplo para ensinar a linguagem.
\end{atencao}

Desde então, esse pequeno programa virou uma espécie de \textbf{tradição} na programação: até hoje, quando alguém aprende uma nova linguagem, é muito comum que o primeiro programa escrito seja, justamente, um ``Hello, World!'' --- seja em C, Python, Java, JavaScript ou qualquer outra linguagem.

\section{\textcolor{azulPrincipal}{Tradutor x Compilador: como o computador ``entende'' o nosso código?}}

Quando escrevemos um programa em C, escrevemos em um arquivo de texto com a extensão \textbf{\texttt{.c}}. Porém, o computador não entende diretamente o que escrevemos --- ele só entende \textbf{código de máquina} (sequências de 0s e 1s).

Para que o nosso código vire algo que o computador consiga executar, precisamos de um \textbf{tradutor}. Existem dois tipos principais de ``tradutor'': o \textbf{compilador} e o \textbf{interpretador}. É comum usarmos a palavra ``tradutor'' como um termo geral para os dois, então vamos entender a diferença:

\subsection*{\textcolor{verdeC}{Compilador}}

O \textbf{compilador} lê \textbf{todo o código de uma vez}, traduz tudo para código de máquina e gera um \textbf{arquivo executável} (por exemplo, \texttt{programa.exe} no Windows). Depois disso, basta executar esse arquivo --- o computador não precisa mais ``olhar'' para o código original.

\begin{itemize}
    \item É o que a linguagem \textbf{C} utiliza.
    \item \textcolor{verdeC}{Vantagem}: o programa final costuma rodar mais rápido, pois já está totalmente traduzido.
    \item \textcolor{laranja}{Analogia}: é como \textbf{traduzir um livro inteiro} antes de entregá-lo ao leitor. O trabalho de tradução é feito uma vez só, no início.
\end{itemize}

\subsection*{\textcolor{verdeC}{Interpretador}}

O \textbf{interpretador} lê o código \textbf{linha por linha}, traduzindo e executando cada linha na hora, sem gerar um arquivo executável separado.

\begin{itemize}
    \item É o que linguagens como \textbf{Python} e \textbf{JavaScript}, por exemplo, costumam utilizar (pelo menos em sua forma mais simples).
    \item \textcolor{verdeC}{Vantagem}: é mais rápido para testar pequenas alterações, já que não é preciso ``recompilar'' tudo.
    \item \textcolor{laranja}{Analogia}: é como ter um \textbf{intérprete simultâneo} em uma palestra, traduzindo frase por frase, em tempo real, conforme o palestrante fala.
\end{itemize}

\subsection*{\textcolor{verdeC}{Resumindo}}

O caminho do nosso código em C é:

\begin{center}
    \fbox{\texttt{Código em C (.c)}} $\rightarrow$ \fbox{\textcolor{verdeC}{Compilador}} $\rightarrow$ \fbox{\texttt{Programa executável (.exe)}}
\end{center}

Ou seja:
\begin{enumerate}
    \item Escrevemos o código-fonte (\texttt{programa.c}).
    \item O compilador lê esse código e verifica se está tudo certo (sem erros de escrita/sintaxe).
    \item Se estiver tudo certo, o compilador gera um arquivo executável.
    \item Esse arquivo executável é o que de fato roda no computador, podendo ser executado quantas vezes quisermos, sem precisar compilar de novo (a menos que o código mude).
\end{enumerate}

\section{\textcolor{azulPrincipal}{Onde escrever e executar nosso código? Conhecendo as IDEs}}

Agora que já sabemos que precisamos de um compilador para transformar nosso código \texttt{.c} em um programa executável, surge a pergunta: \textbf{onde vamos escrever esse código?}

Tecnicamente, dá para escrever um programa em C em qualquer editor de texto simples (como o Bloco de Notas) e depois compilar usando comandos no terminal. Só que, na prática, isso seria bem trabalhoso e cheio de detalhes para quem está começando.

Por isso, usamos uma \textbf{IDE} --- sigla para \textbf{Integrated Development Environment} (em português, \textbf{Ambiente de Desenvolvimento Integrado}). Uma IDE é um programa que reúne, em um só lugar:

\begin{destacado}
\begin{itemize}
    \item Um \textbf{editor de texto} preparado para programação (com cores destacando palavras-chave, identificação de erros, etc.).
    \item Ferramentas para \textbf{compilar} e \textbf{executar} o código com poucos cliques (ou atalhos de teclado).
    \item Recursos extras, como sugestões automáticas de código, organização de arquivos do projeto, depuração (encontrar erros passo a passo), entre outros.
\end{itemize}
\end{destacado}

Pense na IDE como uma \textbf{``oficina completa''}: em vez de você precisar ir buscar cada ferramenta separadamente (editor de texto aqui, compilador ali, terminal em outro lugar), a IDE já entrega tudo organizado em um só ambiente.

\subsection*{\textcolor{verdeC}{Sugestão de IDE: Visual Studio Code (VS Code)}}

Nesta apostila, vamos sugerir o uso do \textbf{Visual Studio Code}, conhecido como \textbf{VS Code}. Ele é:

\begin{itemize}
    \item \textcolor{verdeC}{\faCheckCircle} \textbf{Gratuito} e disponível para Windows, Linux e macOS.
    \item \textcolor{verdeC}{\faCheckCircle} \textbf{Leve} e bastante popular no mercado, usado tanto por iniciantes quanto por profissionais experientes.
    \item \textcolor{verdeC}{\faCheckCircle} \textbf{Extensível}: por meio de extensões (complementos), conseguimos adicionar suporte para a linguagem C, com realce de sintaxe, sugestões de código e integração com o compilador.
\end{itemize}

Para programar em C no VS Code, normalmente também é necessário instalar um \textbf{compilador} no computador (por exemplo, o \textbf{GCC}, um dos compiladores de C mais usados), além de configurar a extensão de C/C++. Esse processo de instalação será detalhado em um material à parte, com o passo a passo de configuração do ambiente.

\section{\textcolor{azulPrincipal}{E se eu ainda não quiser instalar nada? Compiladores online}}

Sabemos que, às vezes, instalar e configurar uma IDE e um compilador pode levar um tempinho, ou você pode estar usando um computador onde não tem permissão para instalar programas (como em alguns laboratórios de escola/faculdade).

\begin{dica}[Compiladores online]
Sites que permitem \textbf{escrever, compilar e executar código C diretamente no navegador}, sem precisar instalar nada no computador.
\end{dica}

Algumas opções bastante usadas:

\begin{itemize}
    \item \textbf{Replit} (\url{https://replit.com})
    \item \textbf{OnlineGDB} (\url{https://www.onlinegdb.com})
    \item \textbf{Programiz Online C Compiler} (\url{https://www.programiz.com/c-programming/online-compiler/})
\end{itemize}

Essas ferramentas são ótimas para:

\begin{itemize}
    \item Testar rapidamente um trecho de código.
    \item Praticar exercícios da apostila sem precisar configurar nada.
    \item Usar em computadores compartilhados, como os de escolas e faculdades.
\end{itemize}

Por outro lado, conforme você for avançando nos estudos e começar a desenvolver projetos maiores, com vários arquivos, é recomendável migrar para uma IDE instalada no seu computador, como o \textbf{VS Code}, que oferece mais recursos e controle sobre o seu ambiente de desenvolvimento.

\begin{dica}[Dica]
Não tem problema nenhum começar pelos compiladores online! O importante é focar em entender a lógica e a sintaxe da linguagem C. Quando se sentir confortável, configure o VS Code com calma, seguindo um tutorial de instalação.
\end{dica}

\section{\textcolor{azulPrincipal}{A estrutura básica de um programa em C}}

Todo programa em C, por mais simples que seja, costuma seguir uma estrutura parecida. Vamos ver o exemplo mais clássico de toda a programação: o programa \textbf{``Olá, Mundo!''}.

\begin{lstlisting}[style=codigoC]
#include <stdio.h>

int main() {

    printf("Ola, Mundo!");

    return 0;
}
\end{lstlisting}

Vamos entender, parte por parte, o que cada linha desse código significa:

\subsection*{\textcolor{verdeC}{\texttt{\#include <stdio.h>}}}
Essa linha avisa ao compilador que vamos usar funções \textbf{prontas} da biblioteca de entrada e saída padrão (do inglês \textit{standard input/output}). É dessa biblioteca que vem, por exemplo, a função \texttt{printf}, que usamos para exibir mensagens na tela.

Pense nisso como ``pegar emprestada uma caixa de ferramentas'' antes de começar o trabalho.

\subsection*{\textcolor{verdeC}{\texttt{int main() \{ ... \}}}}
Toda execução de um programa em C começa pela função \texttt{main}. É como se fosse a ``porta de entrada'' do programa: o computador sempre começa a executar por aqui.

Tudo o que está dentro das chaves \texttt{\{ \}} faz parte do que o programa vai executar.

\subsection*{\textcolor{verdeC}{\texttt{printf("Ola, Mundo!");}}}
A função \texttt{printf} serve para \textbf{exibir algo na tela} (imprimir, mostrar). Nesse caso, ela vai mostrar o texto \texttt{Ola, Mundo!} para o usuário.

Repare que:
\begin{itemize}
    \item O texto que queremos exibir fica entre \textbf{aspas duplas} \texttt{" "}.
    \item A linha termina com \textbf{ponto e vírgula} \texttt{;}. Em C, quase toda instrução termina com ponto e vírgula --- é como o ``ponto final'' de uma frase.
\end{itemize}

\subsection*{\textcolor{verdeC}{\texttt{return 0;}}}
Essa linha indica que o programa terminou \textbf{sem erros}. O número \texttt{0} é uma forma de ``avisar'' ao sistema operacional que tudo correu bem.

\section{\textcolor{azulPrincipal}{Resumindo com uma analogia}}

\begin{dica}[Analogia]
Imagine que escrever um programa em C é como \textbf{escrever uma receita de bolo} para alguém que vai seguir cada passo à risca:

\begin{itemize}
    \item O \textbf{\texttt{\#include}} é como separar os utensílios e ingredientes que vamos usar.
    \item A \textbf{\texttt{main()}} é o título da receita, indicando ``é aqui que a receita começa''.
    \item Cada instrução dentro das chaves é \textbf{um passo da receita}, executado um de cada vez, na ordem em que está escrito.
    \item O \textbf{\texttt{return 0}} é o ``prato pronto, sem nenhum problema na cozinha''.
\end{itemize}
\end{dica}

Com o tempo, essa receita vai ficando mais complexa: vamos aprender a guardar informações (variáveis), tomar decisões (estruturas de controle, como \texttt{if} e \texttt{else}) e repetir passos (estruturas de repetição, como \texttt{for} e \texttt{while}). Mas tudo sempre vai seguir essa mesma base que vimos aqui.

\section{\textcolor{azulPrincipal}{O que vem por aí}}

Nos próximos capítulos, vamos explorar:

\begin{itemize}
    \item \textbf{\textcolor{laranja}{Variáveis}}: como guardar e usar informações dentro do programa.
    \item \textbf{\textcolor{laranja}{Entrada e saída de dados}}: como pedir informações ao usuário (\texttt{scanf}) e exibir resultados (\texttt{printf}).
    \item \textbf{\textcolor{laranja}{Estruturas de controle}}: como fazer o programa tomar decisões e repetir tarefas.
    \item \textbf{\textcolor{laranja}{Exercícios práticos}}: para você colocar a mão na massa e fixar cada conceito aprendido.
\end{itemize}

\vspace{0.5cm}
\begin{center}
    \textcolor{laranja}{\Large \faRocket}
    
    A programação é uma habilidade que se desenvolve com \textbf{prática}. Não tem problema errar, travar ou não entender de primeira --- isso faz parte do processo de qualquer pessoa que está aprendendo a programar. O importante é ir, passo a passo, testando, errando e corrigindo.
    
    \vspace{0.3cm}
    \textbf{\textcolor{azulPrincipal}{Vamos começar!}} \textcolor{laranja}{$\rightarrow$}
\end{center}